\kefalaio{Πρώτοι αριθμοί και ανάλυση: τα δομικά υλικά}{Κεφάλαιο Α.1.5, «Χαρακτήρες διαιρετότητας – ΜΚΔ – ΕΚΠ – Ανάλυση αριθμού σε γινόμενο πρώτων παραγόντων» (δεύτερο μέρος)}

\begin{anagke}
Ο Θοδωρής, ο μεγάλος αδελφός του Άρη, έχει περάσει να τους δει μετά το σχολείο. Στο λύκειο έχει αρχίσει να παίζει με κρυπτογραφία, το πώς φυλάσσονται οι κωδικοί στο διαδίκτυο, και δεν χάνει ευκαιρία να το πει. «Σκέφτομαι δύο πρώτους αριθμούς», λέει στη Λίνα, που τον κοιτάει με ενδιαφέρον. «Το γινόμενό τους είναι 143. Ποιοι είναι;»

Η Λίνα δοκιμάζει. Δεν διαιρείται με το 2, γιατί τελειώνει σε 3. Το άθροισμα ψηφίων $1+4+3=8$, άρα ούτε με το 3. Ούτε με το 5. Με το 7, $7\cdot 20=140$ και περισσεύουν 3, όχι. Με το 11: $11\cdot 13=143$. «Έντεκα και δεκατρία!» φωνάζει. «Μπράβο», λέει ο Θοδωρής. «Τώρα σκέψου ότι ο αριθμός έχει 600 ψηφία. Ένας υπολογιστής πολλαπλασιάζει τους δύο πρώτους σε ένα κλάσμα του δευτερολέπτου. Για να τους βρει από το γινόμενο, θέλει απίστευτα πολύ χρόνο. Πάνω σε αυτό στηρίζονται οι κωδικοί των τραπεζών».

Η Ζωή ρωτάει κάτι πιο απλό. «Δηλαδή κάθε αριθμός είναι φτιαγμένος από πρώτους, όπως ο τοίχος από τούβλα;» Ναι. Και το να βρίσκεις αυτά τα τούβλα είναι δουλειά με μέθοδο.
\end{anagke}

\begin{lexi}{πρώτος, σύνθετος, παράγοντας}{πρώτος, σύνθετος}
Ο \lexiword{πρώτος} αριθμός δεν λέγεται έτσι επειδή έρχεται πρώτος στη σειρά, αλλά επειδή είναι «πρωταρχικός», πρωτογενής: δεν φτιάχνεται από μικρότερους αριθμούς με πολλαπλασιασμό. Είναι το δομικό υλικό.

Ο \lexiword{σύνθετος} αριθμός είναι αυτός που συντίθεται, δηλαδή φτιάχνεται, από άλλους: το 12 είναι $2\cdot 2\cdot 3$. Ο \lexiword{παράγοντας} είναι ο αριθμός που παίρνει μέρος σε έναν πολλαπλασιασμό και «παράγει» το γινόμενο.
\end{lexi}

\begin{metafrasi}
Το 7 έχει διαιρέτες μόνο το 1 και το 7. & Το 7 είναι πρώτος\\
Το 12 έχει διαιρέτες τους 1, 2, 3, 4, 6, 12. & Το 12 είναι σύνθετος\\
Το 60 σαν γινόμενο πρώτων. & $60=2\cdot 2\cdot 3\cdot 5=2^2\cdot 3\cdot 5$\\
Ο $\alpha$ διαιρεί τον $\beta$. & $\beta=\alpha\cdot\kappa$, με $\kappa$ φυσικό\\
Πολλαπλάσια του 4. & $0,\,4,\,8,\,12,\,16,\,20,\,\dots$\\
\end{metafrasi}

\tmima{Η θεωρία}

\tmimasmall{Πολλαπλάσια και διαιρέτες}

\textbf{Πολλαπλάσια} ενός φυσικού $\alpha$ είναι οι αριθμοί που βγαίνουν αν τον πολλαπλασιάσουμε με φυσικούς: $0,\ \alpha,\ 2\alpha,\ 3\alpha,\ 4\alpha,\ \dots$ \textbf{Διαιρέτες} του $\alpha$ είναι όλοι οι αριθμοί που τον διαιρούν ακριβώς. Κάθε αριθμός $\alpha$ έχει διαιρέτες τουλάχιστον τον 1 και τον εαυτό του. Αν ο $\delta$ διαιρεί έναν αριθμό, τότε ο αριθμός είναι πολλαπλάσιο του $\delta$. Είναι το ίδιο πράγμα, από δύο πλευρές.

\begin{orismos}{πρώτος και σύνθετος}
\textbf{Πρώτος} λέγεται ένας φυσικός αριθμός μεγαλύτερος από το 1 που έχει διαιρέτες μόνο τον 1 και τον εαυτό του. Ένας φυσικός μεγαλύτερος από το 1 που έχει και άλλους διαιρέτες λέγεται \textbf{σύνθετος}. Το 1 δεν είναι ούτε πρώτος ούτε σύνθετος.
\end{orismos}

Οι πρώτοι μέχρι το 50 είναι οι: \blue{2, 3, 5, 7, 11, 13, 17, 19, 23, 29, 31, 37, 41, 43, 47}. Ο 2 είναι ο μοναδικός άρτιος πρώτος, γιατί κάθε άλλος άρτιος διαιρείται και με το 2. Οι πρώτοι είναι άπειροι: το απέδειξε ο Ευκλείδης πριν από 2.300 χρόνια.

\tmimasmall{Το κόσκινο του Ερατοσθένη}

Ο Ερατοσθένης, που έζησε στην Αλεξάνδρεια τον 3ο αιώνα π.Χ., βρήκε μια μέθοδο για να «κοσκινίζουμε» τους πρώτους. Γράφεις τους αριθμούς, σβήνεις το 1, μετά κρατάς το 2 και σβήνεις όλα τα πολλαπλάσιά του, μετά κρατάς το επόμενο που δεν έσβησες, το 3, και σβήνεις τα πολλαπλάσιά του, και ούτω καθεξής. Αυτά που μένουν είναι οι πρώτοι.

\begin{center}
\renewcommand{\arraystretch}{1.25}
\setlength{\tabcolsep}{5pt}
\newcommand{\pr}[1]{\textcolor{gblue}{\bfseries #1}}
\begin{tabular}{cccccccccc}
1 & \pr{2} & \pr{3} & 4 & \pr{5} & 6 & \pr{7} & 8 & 9 & 10\\
\pr{11} & 12 & \pr{13} & 14 & 15 & 16 & \pr{17} & 18 & \pr{19} & 20\\
21 & 22 & \pr{23} & 24 & 25 & 26 & 27 & 28 & \pr{29} & 30\\
\pr{31} & 32 & 33 & 34 & 35 & 36 & \pr{37} & 38 & 39 & 40\\
\pr{41} & 42 & \pr{43} & 44 & 45 & 46 & \pr{47} & 48 & 49 & 50\\
\end{tabular}
\end{center}

\tmimasmall{Ανάλυση σε γινόμενο πρώτων παραγόντων}

Κάθε σύνθετος αριθμός γράφεται ως γινόμενο πρώτων, και \textbf{με έναν μόνο τρόπο} (αν αγνοήσουμε τη σειρά των παραγόντων). Γι' αυτό οι πρώτοι είναι τα δομικά υλικά: κάθε αριθμός έχει το δικό του «σύνολο υλικών».

\begin{kanonas}{ανάλυση σε πρώτους παράγοντες}
\begin{enumerate}[leftmargin=1.6em,itemsep=0.2ex,topsep=2pt]
\item Βρίσκω τον μικρότερο πρώτο που διαιρεί τον αριθμό (δοκιμάζω με τη σειρά 2, 3, 5, 7, 11, \dots με τα κριτήρια διαιρετότητας).
\item Διαιρώ και γράφω τον πρώτο στα δεξιά και το πηλίκο από κάτω.
\item Επαναλαμβάνω με το πηλίκο, ώσπου να φτάσω στο 1.
\item Γράφω το γινόμενο των πρώτων και τους ίδιους παράγοντες τους ενώνω σε δυνάμεις.
\end{enumerate}
\end{kanonas}

Αν ο αριθμός είναι πρώτος, δεν χρειάζεται να δοκιμάσεις όλους τους πρώτους. Δοκιμάζεις με τη σειρά και σταματάς όταν το \textbf{πηλίκο γίνει μικρότερο από τον διαιρέτη}. Αν μέχρι τότε δεν βρήκες διαιρέτη, ο αριθμός είναι πρώτος.

\begin{prosoxi}
\begin{itemize}[leftmargin=1.2em,itemsep=0.3ex,topsep=2pt]
\item Το 1 \textbf{δεν είναι πρώτος}. Αν ήταν, το 6 θα γραφόταν και $2\cdot 3$ και $1\cdot 2\cdot 3$ και $1\cdot 1\cdot 2\cdot 3$, και η ανάλυση δεν θα ήταν μοναδική.
\item Ο 2 είναι πρώτος, και μάλιστα ο μόνος άρτιος. Όχι «όλοι οι πρώτοι είναι περιττοί».
\item Ο αριθμός που μοιάζει πρώτος δεν είναι πάντα. Το 91 φαίνεται πρώτος, αλλά $91=7\cdot 13$.
\end{itemize}
\end{prosoxi}

\tmima{Λυμένα παραδείγματα}

\paradeigma{Πρώτος ή σύνθετος;}
\tixero Οι αριθμοί 91 και 97.\par
\tithelo Να αποφασίσω αν είναι πρώτοι.\par
\possindeo Για το 91: δεν διαιρείται με το 2 (τελειώνει σε 1), ούτε με το 3 ($9+1=10$), ούτε με το 5. Με το 7: $7\cdot 13=91$. Βρήκα διαιρέτη, άρα είναι σύνθετος.\par
Για το 97: ούτε με 2, ούτε με 3 ($9+7=16$), ούτε με 5. Με το 7: $97=7\cdot 13+6$, δεν διαιρείται, και το πηλίκο 13 είναι ακόμα μεγαλύτερο από το 7. Με το 11: $97=11\cdot 8+9$, και το πηλίκο 8 είναι μικρότερο από το 11. Σταματάω.\par
\apantisi Το 91 είναι σύνθετος ($7\cdot 13$), το 97 είναι πρώτος.\par
\elegxos Οι πρώτοι κάτω από το 100 είναι 25 και ο τελευταίος από αυτούς είναι το 97.

\paradeigma{Ανάλυση του 2.520}
\tixero Ο αριθμός 2.520.\par
\tithelo Να τον αναλύσω σε γινόμενο πρώτων παραγόντων.\par
\possindeo Διαιρώ διαδοχικά: $2.520:2=1.260$, $1.260:2=630$, $630:2=315$. Το 315 είναι περιττός, άρα δοκιμάζω με το 3: $315:3=105$, $105:3=35$. Μετά με το 5: $35:5=7$. Τέλος $7:7=1$.\par
Οι παράγοντες είναι τρία δυάρια, δύο τριάρια, ένα πεντάρι και ένα εφτάρι.\par
\apantisi $2.520=2^3\cdot 3^2\cdot 5\cdot 7$.\par
\elegxos $8\cdot 9=72$, $72\cdot 5=360$, $360\cdot 7=2.520$. Σωστά.

\paradeigma{Το μυστικό του Θοδωρή}
\tixero Το 143 είναι το γινόμενο δύο πρώτων.\par
\tithelo Ποιοι είναι αυτοί οι πρώτοι.\par
\possindeo Δοκιμάζω τους πρώτους με τη σειρά. Το 2, το 3, το 5 και το 7 δεν διαιρούν το 143. Με το 11: $143:11=13$, και το 13 είναι πρώτος.\par
\apantisi $143=11\cdot 13$.\par
\elegxos $11\cdot 13=143$. Και οι δύο παράγοντες είναι πρώτοι, άρα αυτή είναι η ανάλυση.

\tmima{Ασκήσεις}

\zoni{Το κατάλαβα}{εφαρμογή του κανόνα}

\ex Γράψε όλους τους διαιρέτες των αριθμών 12, 18, 30, 36 και 45.

\ex Γράψε τα έξι πρώτα πολλαπλάσια (μαζί με το 0) των αριθμών 4, 7 και 12.

\ex Ποιοι από τους 1, 2, 9, 13, 21, 29, 33, 41, 51, 57, 61, 87 και 89 είναι πρώτοι; Για τους σύνθετους, γράψε έναν διαιρέτη τους.

\ex Βρες όλους τους πρώτους αριθμούς ανάμεσα στο 50 και στο 80.

\ex Ανάλυσε σε γινόμενο πρώτων παραγόντων (με δυνάμεις): α) 24 \quad β) 60 \quad γ) 72 \quad δ) 100 \quad ε) 150 \quad στ) 504.

\ex Ποιοι αριθμοί έχουν ανάλυση: α) $2^3\cdot 5$ \quad β) $2\cdot 3^2\cdot 7$ \quad γ) $2^2\cdot 3\cdot 5^2$ \quad δ) $3^3\cdot 11$;

\ex Σωστό ή λάθος; α) Κάθε πρώτος αριθμός είναι περιττός. \quad β) Το 1 είναι πρώτος. \quad γ) Το 91 είναι πρώτος. \quad δ) Υπάρχουν άπειροι πρώτοι αριθμοί. \quad ε) Το 0 είναι πολλαπλάσιο κάθε αριθμού. \quad στ) Το 51 είναι πρώτος.

\zoni{Το χρησιμοποιώ}{προβλήματα από τη ζωή}

\ex \textbf{Το μυστικό του Θοδωρή.} Ο Θοδωρής σκέφτεται δύο πρώτους με γινόμενο 221. Ποιοι είναι; Δοκίμασε με τη σειρά τους πρώτους μέχρι να βρεις διαιρέτη.

\ex \textbf{Οι καραμέλες της Μυρτώς.} Η Μυρτώ έχει 48 καραμέλες και θέλει να τις βάλει σε ίσα σακουλάκια, χωρίς να περισσέψει καμία. Πόσες καραμέλες μπορεί να βάλει σε κάθε σακουλάκι; Η Μυρτώ έχει και ένα άλλο βάζο με 47 καραμέλες. Με πόσους τρόπους μπορεί να τις μοιράσει εκεί; Γιατί τόσο λίγους;

\ex \textbf{Οι ομάδες.} Ο γυμναστής θέλει να χωρίσει 60 παιδιά σε ίσες ομάδες. Θέλει τουλάχιστον 2 ομάδες και τουλάχιστον 2 παιδιά σε κάθε ομάδα. Με πόσους διαφορετικούς τρόπους γίνεται; Αν τα παιδιά ήταν 61, πόσοι τρόποι θα υπήρχαν;

\ex \textbf{Η φανέλα του Άρη.} Ο αριθμός στη φανέλα του Άρη είναι πρώτος αριθμός ανάμεσα στο 20 και στο 30, και το άθροισμα των ψηφίων του είναι 11. Ποιος είναι ο αριθμός;

\ex \textbf{Το κόσκινο.} Φτιάξε το κόσκινο του Ερατοσθένη για τους αριθμούς από το 1 ως το 60. Πόσοι πρώτοι υπάρχουν; Πόσους αριθμούς σβήνει το 2, και πόσους ακόμα το 3;

\zoni{Το σκέφτομαι}{μια ερώτηση που πάει πιο μακριά}

\ex Βρες τους διαιρέτες του 6 και του 28, εκτός από τον ίδιο τον αριθμό, και πρόσθεσέ τους. Τι παρατηρείς; Οι μαθηματικοί λένε αυτούς τους αριθμούς «τέλειους».

\ex Δύο πρώτοι που διαφέρουν κατά 2 λέγονται δίδυμοι, όπως οι 3 και 5. Βρες όλα τα ζευγάρια διδύμων πρώτων μέχρι το 50.

\ex Ο 72 είναι $2^3\cdot 3^2$. Γράψε όλους τους αριθμούς της μορφής $2^{\alpha}\cdot 3^{\beta}$, όπου $\alpha=0,\,1,2,\,3$ και $\beta=0,\,1,2$. Τι είναι αυτοί οι αριθμοί για τον 72; Πόσοι είναι;

\ex Η Ζωή ρωτάει «γιατί το 1 δεν είναι πρώτος, αφού διαιρείται μόνο με τον εαυτό του και με το 1;» Εξήγησε με την ανάλυση του 6 γιατί, αν το 1 ήταν πρώτος, δεν θα είχε κάθε αριθμός μία μοναδική ανάλυση.

\ex Ο Θοδωρής ισχυρίζεται ότι $10.403=101\cdot 103$. Έλεγξέ το με πολλαπλασιασμό. Πόσο δύσκολο θα ήταν να βρεις τους δύο παράγοντες χωρίς να τους ξέρεις, δοκιμάζοντας τους πρώτους με τη σειρά;
